\documentclass[11pt]{article}
\usepackage[T1]{fontenc}
\usepackage{textcomp}
\usepackage[margin=0.75in]{geometry}
\usepackage{amsmath,amssymb,amsthm}
\usepackage{array,booktabs}
\usepackage{hyperref}
\setlength{\parindent}{0pt}
\newtheorem{theorem}{Theorem}
\theoremstyle{definition}
\newtheorem{definition}{Definition}
\title{Flow Proof Artifact: parallel-co-interior-sum-two-right}
\author{Generated by \texttt{flow doc proof}}
\date{Flow Proof Book}
\begin{document}
\maketitle
Co-interior angles on parallel lines sum to two right angles.\\[0.5em]
\textbf{Source.} euclid — Elements, Book I, Proposition 29\\[0.5em]
\bigskip
\noindent{\large \textbf{Derived fact 1.} \textit{co-interior alpha plus beta equals two right angles} \hfill the Euclidean plane \cdot parallel lines \cdot co-interior angles sum to two right}\par
\smallskip\noindent\textit{Coordinate: the Euclidean plane \cdot parallel lines \cdot co-interior angles sum to two right}\par
\smallskip\noindent\textit{Source: euclid}\par
\smallskip\noindent\textit{Built on: corresponding angles are equal, for parallel lines in on the Euclidean plane, adjacent angles on a line sum to two right, for intersecting lines in on the Euclidean plane}\par
\medskip
\noindent\textbf{Goal.} co-interior alpha plus beta equals two right angles\par
\noindent$\displaystyle \angle \text{co interior alpha} + \angle \text{co interior beta} = 180^\circ$\par
\medskip
\noindent
\begin{tabular}{@{} >{\raggedright\arraybackslash}p{0.06\textwidth} >{\raggedright\arraybackslash}p{0.47\textwidth} >{\raggedleft\arraybackslash}p{0.41\textwidth} @{}}
\textbf{\#} & \textbf{Proof} & \textbf{Mathematics} \\
\midrule
\textbf{1.} & We prove that co-interior angles sum to two right for parallel lines in the Euclidean plane. & \\[0.45em]
\textbf{2.} & We invoke the derived fact governing parallel lines in the Euclidean plane: corresponding angles are equal, for parallel lines in on the Euclidean plane. & \\[0.45em]
\textbf{3.} & We invoke the derived fact governing intersecting lines in the Euclidean plane: adjacent angles on a line sum to two right, for intersecting lines in on the Euclidean plane. & \\[0.45em]
\textbf{4.} & From step 2 and step 3, this implies angle co interior alpha plus angle co interior beta equals two right angles. Hence proven. & $\displaystyle \angle \text{co interior alpha} + \angle \text{co interior beta} = 180^\circ$ \\[0.45em]
\bottomrule
\end{tabular}
\label{thm:Geometry-parallel lines-co_interior_angles_sum_to_two_right}
\par\medskip\hrule
\end{document}
